\exercisesection{正则环}

\begin{enumerate}
\def\labelenumi{\arabic{enumi})}
\item
  设 A 为 Noether 局部环。为使 A 正则，必须且只需 \(\mathrm{di}_{\mathrm{A}}(\kappa_{\mathrm{A}})<+\infty\)。
\item
  设 A 为正则环，p 为 A 的一个素理想。证明 A 关于 p-进拓扑的完备化是整环。
\item
  设 A 为环，\(\mathbf{T} = (T_i)_{i\in I}\) 为一个有限的未定元族。证明下列条件等价：
\end{enumerate}

\begin{enumerate}
\def\labelenumi{(\roman{enumi})}
\item
  环 A 是正则的；
\item
  环 A{[}T{]} 是正则的；
\item
  环 A{[}{[}T{]}{]} 是正则的。
\end{enumerate}

\begin{enumerate}
\def\labelenumi{\arabic{enumi})}
\setcounter{enumi}{3}
\item
  设 A 为维数 \(\leqslant 2\) 的 Noether 环。证明：为使 A 正则，必须且只需每个有限生成 A-模的对偶都是投射的（即设 \(\mathfrak{p}\) 取遍 \(\Lambda\) 的素理想；若第二个条件满足，先证明 \(\mathrm{A}_{\mathfrak{p}}\) 在 \(\operatorname{prof}(A_{\mathfrak{p}}) \leqslant 1\) 时是正则的，再证明 \(\mathrm{dp}_{\mathrm{A}_{\mathfrak{p}}}(\kappa(\mathfrak{p})) = 2\) 在 \(\operatorname{prof}(A_{\mathfrak{p}}) = 2\) 时成立）。
\item
  设 A 为同调维数有限的 Noether 局部环；试由 A, X, 第~208 页习题 16 推出 A 正则这一事实的另一证明（注意有 dim \(_{\kappa_A}\) \(\mathfrak{m}_A/\mathfrak{m}_A^2 \leqslant \mathrm{dh}(A) \leqslant \mathrm{prof}(A)\)）。
\item
  设 A 为正则局部环，\((x_{1},\ldots,x_{n})\) 与 \((y_{1},\ldots,y_{n})\) 为 \(m_{A}\) 中元素的极大割序列，分别生成理想 x 与 y。假设 \(y \subset x\)，于是存在 A 的元素 \(a_{ij}\)，使得 \(y_{i} = \sum_{j} a_{ij} x_{j}\) 对 \(1 \leqslant i \leqslant n\) 成立。证明 \(\det(a_{ij})\) 在 A/y 中的类不依赖于所选取的矩阵 \((a_{ij})\)，且其零化子是 x 的像；特别地，有 \(\det(a_{ij}) \neq 0\)（证明典范同态 \(\operatorname{Ext}_{\mathrm{A}}^{n}(\mathrm{A}/x,\mathrm{A}) \to \operatorname{Ext}_{\mathrm{A}}^{n}(\Lambda/\mathfrak{y},\Lambda)\) 是单射，再利用 Koszul 复形证明它与由比率为 \(\det(a_{ij})\) 的乘法映射诱导的同态 A/x → A/y 等同）。
\end{enumerate}

设 A 为维数 3 的正则局部环，f 为 \(m_{A}\) 的一个非零元素。证明下列条件等价：

\begin{enumerate}
\def\labelenumi{(\roman{enumi})}
\item
  环 A/(f) 不是因子分解整环；
\item
  存在整数 \(n \geqslant 2\) 以及一个矩阵 \(X \in \mathbf{M}_n(\mathrm{A})\)，其元素属于 \(\mathfrak{m}_{\mathrm{A}}\)，使得 \(f = \det(\mathrm{X})\)。
\end{enumerate}

（在条件 (ii) 下，证明 A/(f) 中由 X 的子式 \(X^{11}, \ldots, X^{1n}\)（其中 \(X^{1p}\) 由删去第一行与指标为 p 的列而得）的类生成的理想高度为 1 但不是主理想。在条件 (i) 下，设 p 为 A 的一个高度 2 的素理想，其在 \(\Lambda/(f)\) 中的像不是主理想；由 §3 习题 11 推出 p 由矩阵 \(Y \in \mathbf{M}_{n,n-1}(\mathbf{A})\)（元素属于 \(m_{\Lambda}\)）的 n-1 阶子式生成。把 f 写成矩阵 \(\widetilde{Y}\)（由对 Y 添加一列 \(Y_{n}\) 而得）的行列式；若 \(Y_{n}\) 中有一个元素可逆，则化为 \(Y_{n}\) 属于 \(\Lambda^{n}\) 的典范基的情形。

¶ 8) 设 A 为正则局部环，\(\mathbf{x} = (x_1, \ldots, x_d)\) 为 A 的一个坐标系，\(\rho : A \to B\) 为使 B 成为有限生成 A-模的单射环同态。证明下列条件等价：

\begin{enumerate}
\def\labelenumi{(\roman{enumi})}
\item
  对每个整数 \(p\)，元素 \(\rho(x_1 \ldots x_d)^p\) 不属于由 \(\rho(x_1^{p+1}), \ldots, \rho(x_d^{p+1})\) 生成的 B 的理想；
\item
  A-模 \(\rho(A)\) 是 B 的直和因子。
\end{enumerate}

当 A 是 Q-代数时这些条件成立（§ 2 习题 3）。

为证 (ii)⇒(i)，仿照该处的证明。对逆蕴含，化为 A 为完备的情形，并对 \(p \geqslant 0\) 令 \(A_{p} = A/(x_{1}^{p+1}, \ldots, x_{d}^{p+1})\)，\(B_{p} = B/(x_{1}^{p+1}, \ldots, x_{d}^{p+1})\)。注意 \((x_{1} \ldots x_{d})^{p}\) 的类生成 \(A_{p}\) 的底座，由此推出同态 \(\rho_{p}: A_{p} \to B_{p}\)（由 \(\rho\) 诱导）是单射。证明 \(\rho_{p}\) 容许一个 \(A_{p}\)-线性收缩。最后利用 E, III, 第~60 页，例 II。

\begin{enumerate}
\def\labelenumi{\arabic{enumi})}
\setcounter{enumi}{8}
\tightlist
\item
  设 \(\rho: A \to B\) 为 Noether 局部环之间的一个局部同态，\(N\) 为一个有限生成的 Macaulay B-模。假设环 \(A\) 正则，且我们有 \(\dim_{B}(N) = \dim(A) + \dim_{B/\mathfrak{m}_{A}B}(N/\mathfrak{m}_{A}N)\)。证明 A-模 \(N\) 是平坦的（对 \(\dim(A)\) 作归纳论证；若 \(x \in \mathfrak{m}_{A} - \mathfrak{m}_{A}^{2}\)，将证明 \(x\) 是 \(N\)-正则的，且同态 A/xA → B/xB 与 (B/xB)-模 N/xN 满足与 ρ 和 N 相同的假设）。
\end{enumerate}

¶ 10) 设 A 为 Noether 环，M 为有限生成 A-模。对每个整数 \(n \geqslant 0\)，用 \(X_{n}\) 表示元素 \(\mathfrak{p}\)（属于 \(\operatorname{Spec}(A)\)，满足 \(\dim_{A_{\mathfrak{p}}}(M_{\mathfrak{p}}) - \operatorname{prof}_{A_{\mathfrak{p}}}(M_{\mathfrak{p}}) \geqslant n\)）所成的集合。设 \(k\) 为一个整数 \(\geqslant 0\)。

\begin{enumerate}
\def\labelenumi{\alph{enumi})}
\item
  为使 M 满足性质 \(S_{k}\)（§ 1，习题 8），必须且只需 codim( \(X_{n}\)，Supp(M)) \(\geqslant n + k\) 对每个 \(n \geqslant 0\) 成立。
\item
  假设 \(X_{n}\) 在 \(\operatorname{Spec}(A)\) 中对每个 \(n \geqslant 0\) 都是闭的。设 p 为 \(\operatorname{Supp}(M)\) 的一个元素。为使 \(A_{p}\)-模 \(M_{p}\) 满足性质 \(S_{k}\)（§ 1，习题 8），必须且只需 \(\operatorname{codim}(X_{n}^{\prime}, \operatorname{Supp}(M)) \geqslant n + k\) 对每个 \(n \geqslant 0\) 以及每个分支 \(X_{n}^{\prime}\)（属于 \(X_{n}\) 且包含 p）成立。
\item
  假设环 A 可表现。由 b) 推出：A 中使 \(M_{p}\) 满足性质 \(S_{k}\) 的素理想 p 所成的集合在 Spec(A) 中是开的。
\end{enumerate}

\begin{enumerate}
\def\labelenumi{\arabic{enumi})}
\setcounter{enumi}{10}
\tightlist
\item
  设 A 为环 \(\mathbf{Z}[\mathbf{X}]\)，其中 \(\mathbf{X} = (X_i)_{i\in I}\) 为一个有限的未定元族。设 J 为 A 的一个由单项式生成的理想。
\end{enumerate}

\begin{enumerate}
\def\labelenumi{\alph{enumi})}
\item
  设 \(i \in I\)；设生成 J 的单项式为 \(\mathrm{P}_{1}, \ldots, \mathrm{P}_{r}\) 与 \(\mathrm{Q}_{1}, \ldots, \mathrm{Q}_{s}\)，其中 \(\mathrm{P}_{j}\) 被 \(\mathrm{X}_{i}\) 整除而 \(\mathrm{Q}_{k}\) 则否。证明传递子 J: \(\mathrm{X}_{i}\)（由满足 \(\mathrm{X}_{i} \mathrm{P} \in \mathrm{J}\) 的元素 P 所组成的 A 的理想）由单项式 \(\mathrm{X}_{i}^{-1} \mathrm{P}_{1}, \ldots, \mathrm{X}_{i}^{-1} \mathrm{P}_{r}\) 与 \(\mathrm{Q}_{1}, \ldots, \mathrm{Q}_{s}\) 生成。
\item
  由此推出：若 \(J: X_i\) 等于 \(J\) 或 \(A\) 对每个 \(i \in I\) 成立，则 \(J\) 是变量 \(X_i\)（它们属于 \(J\)）所生成的理想。
\item
  设 M 为一个 A-模，p 为一个整数。假设 \(\operatorname{Tor}_{p}^{A}(A/J',M)\)（相应地，\(\operatorname{Ext}_{A}^{p}(A/J',M)\)）对每个理想 \(J'\)（包含 J 且由某些 \(X_{i}\) 生成）均为零。证明 \(\operatorname{Tor}_{p}^{A}(A/J,M)\)（相应地，\(\operatorname{Ext}_{A}^{p}(A/J,M)\)）为零（取一个理想 \(J'\)，它包含 J、由单项式生成、使 \(\operatorname{Tor}_{p}^{A}(A/J',M) \neq 0\) 且对这些性质为极大；由正合列
\[
0 \rightarrow \mathrm{A} / (\mathrm{J} ^ {\prime}: \mathrm{X} _ {i}) \longrightarrow \mathrm{A} / \mathrm{J} ^ {\prime} \longrightarrow \mathrm{A} / (\mathrm{J} ^ {\prime} + \mathrm{AX} _ {i}) \rightarrow 0
\]

以及 b) 推出 J' 由它所包含的元素 \(X_{i}\) 生成，这与假设矛盾）。

\item
  由此推出不等式 \(\mathrm{dp}_{\mathrm{A}}(\mathrm{A}/\mathrm{J}) \leqslant \mathrm{Card}(\mathrm{I})\)。

\item
  设 A 为环，\(\mathbf{x} = (x_{i})_{i \in I}\) 为 A 中元素的一个有限族；假设对 I 的每个子集 I′，族 \((x_{i})_{i \in I'}\) 对 A 是完全割的。设 J 为由元素 \(x^{\alpha}\) 生成的 A 的理想，其中 \(\alpha\) 取遍 \(N^{1}\) 的一个有限子集 F。证明不等式 dp \(_{A}\)(A/J) ≤ Card(I)（设 J \(_{0}\) 为环 A \(_{0}\) = Z{[}(X \(_{i}\)) \(_{i \in I}\){]} 中由诸 \(X^{\alpha}\)（\(\alpha \in F\)）生成的理想；由 c) 推出 Tor \(_{p}^{A_{0}}\)(A \(_{0}\)/J \(_{0}\), A) 为零，并考虑 A \(_{0}\)-模 A \(_{0}\)/J \(_{0}\) 的一个长度 ≤ n 的投射分解来完成证明）。
\end{enumerate}

¶ 12) 设 A 为正则局部环，M 为有限生成 A-模，且 A-模 End \(_{A}\)(M) 是自由的。我们旨在证明下列条件等价：

\begin{enumerate}
\def\labelenumi{(\roman{enumi})}
\item
  A-模 M 是自由的；
\item
  A-模 M 是自反的；
\item
  \(\operatorname{Ext}_{\mathrm{A}}^{1}(\mathbf{M},\mathbf{M}) = 0\)
\end{enumerate}

\begin{enumerate}
\def\labelenumi{\alph{enumi})}
\item
  证明蕴含关系 (iii) \(\Rightarrow\) (ii)（利用 VII, § 4 习题 2）。
\item
  证明逆蕴含（当 dim(A) ≤ 2 时利用习题 4，然后利用 § 1 习题 18 对 dim(A) 作归纳论证）。
\item
  假设 \(\dim(A) \leqslant 3\)。证明蕴含关系 (ii) \(\Rightarrow\) (i)（当 \(\dim(A) \leqslant 2\) 时利用习题 4；若 \(\dim(A) = 3\)，注意 (ii) 蕴含 \(\mathrm{dp}_{A}(M) \leqslant 1\)，再注意 (iii) 与 \(n^{\circ} 3\) 的注 2 蕴含 M 是投射的）。
\item
  假设 M 是自反的。设 x 为 \(m_{A} = m_{A}^{2}\) 的一个元素；用 B 表示环 A/xA，用 N 表示 B-模 M/xM，用 \(N^{*}\) 表示它的对偶。利用 b)，证明 B-模 \(\operatorname{End}_{\mathbb{B}}(\mathbb{N})\) 是自由的，再证明 \(\operatorname{End}_{\mathbb{B}}(\mathbb{N}^{*})\) 是自由的（由 VII, § 4, 第 2 节推出 \(\operatorname{End}_{\mathbb{B}}(\mathbb{N}^{*})\) 与 \(\operatorname{End}_{\mathbb{B}}(\mathbb{N})\) 的双对偶等同）。
\item
  证明 (ii) 蕴含 (i)（对 A 的维数作归纳论证，设其 ≥ 4；采用 d) 的记号，归纳假设蕴含 B-模 N* 是自由的；由此并利用 § 1 习题 18 与第 4 节命题 7，推出 prof \(_{A}\)(Ext \(^{1}_{A}\)(M, B)) ≥ 1，再推出按归纳假设为有限长度的 Ext \(^{1}_{A}\)(M, A) 为零；断言 B-模 M*/xM* 与 N* 等同，从而是自由的，进而 A-模 M* 是自由的，最后 M 是自由的。）
\item
  由这些结果推出环 A 是因子分解整环这一事实的另一证明。
\end{enumerate}

\begin{enumerate}
\def\labelenumi{\arabic{enumi})}
\setcounter{enumi}{11}
\tightlist
\item
  设 A 为环，\((x_{1},\ldots,x_{n})\) 为 A 中元素的一个族，生成理想 I。若族 \((x_{1},\ldots,x_{n})\) 的诸 \(x_{i}\) 在 I/I \(^{2}\) 中的类构成该 A/l-模的一个基，则称该族是 A-无关的。
\end{enumerate}

\begin{enumerate}
\def\labelenumi{\alph{enumi})}
\item
  每个对 A 的完全割族都是 A-无关的。
\item
  设 \((x_{1},\ldots,x_{n})\) 为一个 A-无关族，且 \(x_{1}=x_{1}^{\prime}x_{1}^{\prime\prime}\)。证明族 \((x_{1}^{\prime},\ldots,x_{n})\) 与 \((x_{1}^{\prime\prime},\ldots,x_{n})\) 都是 A-无关的，并且有
\end{enumerate}

\[
\lg_ {\mathrm{A}} \left(\mathrm{A} / \left(x _ {1}, \dots , x _ {n}\right)\right) = \lg_ {\mathrm{A}} \left(\mathrm{A} / \left(x _ {1} ^ {\prime}, \dots , x _ {n}\right)\right) + \lg_ {\mathrm{A}} \left(\mathrm{A} / \left(x _ {1} ^ {\prime \prime}, \dots , x _ {n}\right)\right)
\]

（注意商 \((x_{1}^{\prime},\ldots,x_{n})/(x_{1},\ldots,x_{n})\) 同构于 \(A/(x_{1}^{\prime\prime},\ldots,x_{n})\)）。

¶ 13) 设 \(p\) 为素数，\(q\) 为 \(p\) 的一个幂。设 A 为特征 \(p\) 的约化 Noether 局部环。用 \(\mathrm{A}^q\) 表示由诸元素 \(a^q\)（\(a \in \mathrm{A}\)）组成的 A 的子环。证明：为使 A 正则，必须且只需 \(\mathrm{A}^q\)-模 A 是平坦的。

（化为 A 为完备的情形，于是 A 是形式幂级数代数 \(B = \kappa_{A}[[X_{1}, \ldots, X_{n}]]\) 对某个理想 b 的商，该理想可假设含于 \(m_{B}^{2}\)。若 A 正则，则 b = 0。若 A 在 \(A^{q}\) 上平坦，设 \(x_{i}\) 为 \(X_{i}\) 在 A 中的像；证明族 \((x_{1}^{q}, \ldots, x_{n}^{q})\) 是 A-无关的（习题 12），由此推出等式 \(\lg_{\mathrm{A}}(\mathrm{A}/(x_{1}^{q}, \ldots, x_{n}^{q})) = q^{n}\)。证明这蕴含 \(b \subset (X_{1}^{q}, \ldots, X_{n}^{q})\)；类似地证明对每个 s 有 \(b \subset (X_{1}^{q^{s}}, \ldots, X_{n}^{q^{s}})\)，从而最终 b = 0。）

